% Bewertungslogik: Defizite D1-D4 -> Forschungsfragen F1-F4 -> Anforderungen A1-A7 -> Versuche V1-V11
% -> Abschnitte 8.1-8.4 -> Abgleich (Abschnitt 8.5.1). Eigene Darstellung, keine Fremdquelle.
% Zuordnung nach review/plan-kap5.md Abschnitt 6 und 7.2 (Datenbasis: Planung, kein Messergebnis).
% Haupt-Zuordnungen sind gezeichnet; V2 stützt zusätzlich A1 (K1.1) und A5 (K5.6) und steht in der Caption.
% Kennzahlen K stehen nicht einzeln in der Grafik, sondern in tab:kennzahlen.
% Anforderungen in der Reihenfolge der Fragen angeordnet (A2, A3, A4, A5, A6, A1, A7), damit sich keine Linien kreuzen.
% Stil: ba-grafik.tex. Einbinden: % Bewertungslogik: Defizite D1-D4 -> Forschungsfragen F1-F4 -> Anforderungen A1-A7 -> Versuche V1-V11
% -> Abschnitte 8.1-8.4 -> Abgleich (Abschnitt 8.5.1). Eigene Darstellung, keine Fremdquelle.
% Zuordnung nach review/plan-kap5.md Abschnitt 6 und 7.2 (Datenbasis: Planung, kein Messergebnis).
% Haupt-Zuordnungen sind gezeichnet; V2 stützt zusätzlich A1 (K1.1) und A5 (K5.6) und steht in der Caption.
% Kennzahlen K stehen nicht einzeln in der Grafik, sondern in tab:kennzahlen.
% Anforderungen in der Reihenfolge der Fragen angeordnet (A2, A3, A4, A5, A6, A1, A7), damit sich keine Linien kreuzen.
% Stil: ba-grafik.tex. Einbinden: % Bewertungslogik: Defizite D1-D4 -> Forschungsfragen F1-F4 -> Anforderungen A1-A7 -> Versuche V1-V11
% -> Abschnitte 8.1-8.4 -> Abgleich (Abschnitt 8.5.1). Eigene Darstellung, keine Fremdquelle.
% Zuordnung nach review/plan-kap5.md Abschnitt 6 und 7.2 (Datenbasis: Planung, kein Messergebnis).
% Haupt-Zuordnungen sind gezeichnet; V2 stützt zusätzlich A1 (K1.1) und A5 (K5.6) und steht in der Caption.
% Kennzahlen K stehen nicht einzeln in der Grafik, sondern in tab:kennzahlen.
% Anforderungen in der Reihenfolge der Fragen angeordnet (A2, A3, A4, A5, A6, A1, A7), damit sich keine Linien kreuzen.
% Stil: ba-grafik.tex. Einbinden: % Bewertungslogik: Defizite D1-D4 -> Forschungsfragen F1-F4 -> Anforderungen A1-A7 -> Versuche V1-V11
% -> Abschnitte 8.1-8.4 -> Abgleich (Abschnitt 8.5.1). Eigene Darstellung, keine Fremdquelle.
% Zuordnung nach review/plan-kap5.md Abschnitt 6 und 7.2 (Datenbasis: Planung, kein Messergebnis).
% Haupt-Zuordnungen sind gezeichnet; V2 stützt zusätzlich A1 (K1.1) und A5 (K5.6) und steht in der Caption.
% Kennzahlen K stehen nicht einzeln in der Grafik, sondern in tab:kennzahlen.
% Anforderungen in der Reihenfolge der Fragen angeordnet (A2, A3, A4, A5, A6, A1, A7), damit sich keine Linien kreuzen.
% Stil: ba-grafik.tex. Einbinden: \input{Bilder/tikz/bewertungslogik} in einer figure-Umgebung. Nicht skalieren.
\begin{tikzpicture}[ba]
  \tikzset{
    klein/.style={ba block, minimum width=12mm, minimum height=7mm, inner sep=2pt},
    frage/.style={ba block hervor, minimum width=12mm, minimum height=7mm, inner sep=2pt},
    vers/.style={ba block, minimum width=21mm, inner sep=2pt},
    absch/.style={ba block, minimum width=21mm, inner xsep=5pt, inner ysep=2pt},
    kopf/.style={font=\footnotesize, text=baGrau, anchor=base},
  }
  % ---------- Spaltenköpfe ----------
  \node[kopf] at (0.6,0.65)  {Defizite};
  \node[kopf] at (2.9,0.65)  {Fragen};
  \node[kopf] at (5.2,0.65)  {Anforderungen};
  \node[kopf] at (8.2,0.65)  {Versuche};
  \node[kopf] at (11.3,0.65) {Abschnitte};
  \node[kopf] at (14.4,0.65) {Abgleich};
  % ---------- Anforderungen (Reihenfolge nach Fragen) ----------
  \node[klein] (A2) at (5.2,0)     {A2};
  \node[klein] (A3) at (5.2,-0.95) {A3};
  \node[klein] (A4) at (5.2,-1.9)  {A4};
  \node[klein] (A5) at (5.2,-2.85) {A5};
  \node[klein] (A6) at (5.2,-3.8)  {A6};
  \node[klein] (A1) at (5.2,-4.75) {A1};
  \node[klein] (A7) at (5.2,-5.7)  {A7};
  % ---------- Fragen und Defizite ----------
  \node[frage] (F1) at (2.9,0)      {F1};
  \node[frage] (F2) at (2.9,-1.425) {F2};
  \node[frage] (F3) at (2.9,-3.325) {F3};
  \node[frage] (F4) at (2.9,-5.225) {F4};
  \node[klein] (D1) at (0.6,0)      {D1};
  \node[klein] (D2) at (0.6,-1.425) {D2};
  \node[klein] (D3) at (0.6,-3.325) {D3};
  \node[klein] (D4) at (0.6,-5.225) {D4};
  % ---------- Versuche und Abschnitte (nach Abschnitten gruppiert) ----------
  \node[vers] (VA) at (8.2,0)      {V1 bis V3};
  \node[vers, minimum height=16.5mm] (VB) at (8.2,-1.425) {V4 bis V7};
  \node[vers, minimum height=16.5mm] (VC) at (8.2,-3.325) {V8, V9};
  \node[vers, minimum height=16.5mm] (VD) at (8.2,-5.225) {V10, V11};
  \node[absch] (SA) at (11.3,0)      {Abschnitt~\ref{sec:eval_f1}};
  \node[absch, minimum height=16.5mm] (SB) at (11.3,-1.425) {Abschnitt~\ref{sec:eval_f2}};
  \node[absch, minimum height=16.5mm] (SC) at (11.3,-3.325) {Abschnitt~\ref{sec:eval_f3}};
  \node[absch, minimum height=16.5mm] (SD) at (11.3,-5.225) {Abschnitt~\ref{sec:eval_f4}};
  \node[ba block, minimum width=22mm, minimum height=64mm, text width=20mm] (AB) at (14.4,-2.85)
    {Abgleich\\Abschnitt\\\ref{sec:abgleich}};
  % ---------- Verbindungen ----------
  \foreach \a/\b in {D1/F1, D2/F2, D3/F3, D4/F4} \draw[ba pfeil] (\a) -- (\b);
  \draw[ba pfeil] (F1.east) -- (A2.west);
  \draw[ba pfeil] (F2.east) -- (A3.west);
  \draw[ba pfeil] (F2.east) -- (A4.west);
  \draw[ba pfeil] (F3.east) -- (A5.west);
  \draw[ba pfeil] (F3.east) -- (A6.west);
  \draw[ba pfeil] (F4.east) -- (A1.west);
  \draw[ba pfeil] (F4.east) -- (A7.west);
  \draw[ba pfeil] (A2.east) -- (A2 -| VA.west);
  \draw[ba pfeil] (A3.east) -- (A3 -| VB.west);
  \draw[ba pfeil] (A4.east) -- (A4 -| VB.west);
  \draw[ba pfeil] (A5.east) -- (A5 -| VC.west);
  \draw[ba pfeil] (A6.east) -- (A6 -| VC.west);
  \draw[ba pfeil] (A1.east) -- (A1 -| VD.west);
  \draw[ba pfeil] (A7.east) -- (A7 -| VD.west);
  \foreach \v/\s in {VA/SA, VB/SB, VC/SC, VD/SD} \draw[ba pfeil] (\v) -- (\s);
  \draw[ba pfeil] (SA.east) -- (SA -| AB.west);
  \draw[ba pfeil] (SB.east) -- (SB -| AB.west);
  \draw[ba pfeil] (SC.east) -- (SC -| AB.west);
  \draw[ba pfeil] (SD.east) -- (SD -| AB.west);
\end{tikzpicture}
 in einer figure-Umgebung. Nicht skalieren.
\begin{tikzpicture}[ba]
  \tikzset{
    klein/.style={ba block, minimum width=12mm, minimum height=7mm, inner sep=2pt},
    frage/.style={ba block hervor, minimum width=12mm, minimum height=7mm, inner sep=2pt},
    vers/.style={ba block, minimum width=21mm, inner sep=2pt},
    absch/.style={ba block, minimum width=21mm, inner xsep=5pt, inner ysep=2pt},
    kopf/.style={font=\footnotesize, text=baGrau, anchor=base},
  }
  % ---------- Spaltenköpfe ----------
  \node[kopf] at (0.6,0.65)  {Defizite};
  \node[kopf] at (2.9,0.65)  {Fragen};
  \node[kopf] at (5.2,0.65)  {Anforderungen};
  \node[kopf] at (8.2,0.65)  {Versuche};
  \node[kopf] at (11.3,0.65) {Abschnitte};
  \node[kopf] at (14.4,0.65) {Abgleich};
  % ---------- Anforderungen (Reihenfolge nach Fragen) ----------
  \node[klein] (A2) at (5.2,0)     {A2};
  \node[klein] (A3) at (5.2,-0.95) {A3};
  \node[klein] (A4) at (5.2,-1.9)  {A4};
  \node[klein] (A5) at (5.2,-2.85) {A5};
  \node[klein] (A6) at (5.2,-3.8)  {A6};
  \node[klein] (A1) at (5.2,-4.75) {A1};
  \node[klein] (A7) at (5.2,-5.7)  {A7};
  % ---------- Fragen und Defizite ----------
  \node[frage] (F1) at (2.9,0)      {F1};
  \node[frage] (F2) at (2.9,-1.425) {F2};
  \node[frage] (F3) at (2.9,-3.325) {F3};
  \node[frage] (F4) at (2.9,-5.225) {F4};
  \node[klein] (D1) at (0.6,0)      {D1};
  \node[klein] (D2) at (0.6,-1.425) {D2};
  \node[klein] (D3) at (0.6,-3.325) {D3};
  \node[klein] (D4) at (0.6,-5.225) {D4};
  % ---------- Versuche und Abschnitte (nach Abschnitten gruppiert) ----------
  \node[vers] (VA) at (8.2,0)      {V1 bis V3};
  \node[vers, minimum height=16.5mm] (VB) at (8.2,-1.425) {V4 bis V7};
  \node[vers, minimum height=16.5mm] (VC) at (8.2,-3.325) {V8, V9};
  \node[vers, minimum height=16.5mm] (VD) at (8.2,-5.225) {V10, V11};
  \node[absch] (SA) at (11.3,0)      {Abschnitt~\ref{sec:eval_f1}};
  \node[absch, minimum height=16.5mm] (SB) at (11.3,-1.425) {Abschnitt~\ref{sec:eval_f2}};
  \node[absch, minimum height=16.5mm] (SC) at (11.3,-3.325) {Abschnitt~\ref{sec:eval_f3}};
  \node[absch, minimum height=16.5mm] (SD) at (11.3,-5.225) {Abschnitt~\ref{sec:eval_f4}};
  \node[ba block, minimum width=22mm, minimum height=64mm, text width=20mm] (AB) at (14.4,-2.85)
    {Abgleich\\Abschnitt\\\ref{sec:abgleich}};
  % ---------- Verbindungen ----------
  \foreach \a/\b in {D1/F1, D2/F2, D3/F3, D4/F4} \draw[ba pfeil] (\a) -- (\b);
  \draw[ba pfeil] (F1.east) -- (A2.west);
  \draw[ba pfeil] (F2.east) -- (A3.west);
  \draw[ba pfeil] (F2.east) -- (A4.west);
  \draw[ba pfeil] (F3.east) -- (A5.west);
  \draw[ba pfeil] (F3.east) -- (A6.west);
  \draw[ba pfeil] (F4.east) -- (A1.west);
  \draw[ba pfeil] (F4.east) -- (A7.west);
  \draw[ba pfeil] (A2.east) -- (A2 -| VA.west);
  \draw[ba pfeil] (A3.east) -- (A3 -| VB.west);
  \draw[ba pfeil] (A4.east) -- (A4 -| VB.west);
  \draw[ba pfeil] (A5.east) -- (A5 -| VC.west);
  \draw[ba pfeil] (A6.east) -- (A6 -| VC.west);
  \draw[ba pfeil] (A1.east) -- (A1 -| VD.west);
  \draw[ba pfeil] (A7.east) -- (A7 -| VD.west);
  \foreach \v/\s in {VA/SA, VB/SB, VC/SC, VD/SD} \draw[ba pfeil] (\v) -- (\s);
  \draw[ba pfeil] (SA.east) -- (SA -| AB.west);
  \draw[ba pfeil] (SB.east) -- (SB -| AB.west);
  \draw[ba pfeil] (SC.east) -- (SC -| AB.west);
  \draw[ba pfeil] (SD.east) -- (SD -| AB.west);
\end{tikzpicture}
 in einer figure-Umgebung. Nicht skalieren.
\begin{tikzpicture}[ba]
  \tikzset{
    klein/.style={ba block, minimum width=12mm, minimum height=7mm, inner sep=2pt},
    frage/.style={ba block hervor, minimum width=12mm, minimum height=7mm, inner sep=2pt},
    vers/.style={ba block, minimum width=21mm, inner sep=2pt},
    absch/.style={ba block, minimum width=21mm, inner xsep=5pt, inner ysep=2pt},
    kopf/.style={font=\footnotesize, text=baGrau, anchor=base},
  }
  % ---------- Spaltenköpfe ----------
  \node[kopf] at (0.6,0.65)  {Defizite};
  \node[kopf] at (2.9,0.65)  {Fragen};
  \node[kopf] at (5.2,0.65)  {Anforderungen};
  \node[kopf] at (8.2,0.65)  {Versuche};
  \node[kopf] at (11.3,0.65) {Abschnitte};
  \node[kopf] at (14.4,0.65) {Abgleich};
  % ---------- Anforderungen (Reihenfolge nach Fragen) ----------
  \node[klein] (A2) at (5.2,0)     {A2};
  \node[klein] (A3) at (5.2,-0.95) {A3};
  \node[klein] (A4) at (5.2,-1.9)  {A4};
  \node[klein] (A5) at (5.2,-2.85) {A5};
  \node[klein] (A6) at (5.2,-3.8)  {A6};
  \node[klein] (A1) at (5.2,-4.75) {A1};
  \node[klein] (A7) at (5.2,-5.7)  {A7};
  % ---------- Fragen und Defizite ----------
  \node[frage] (F1) at (2.9,0)      {F1};
  \node[frage] (F2) at (2.9,-1.425) {F2};
  \node[frage] (F3) at (2.9,-3.325) {F3};
  \node[frage] (F4) at (2.9,-5.225) {F4};
  \node[klein] (D1) at (0.6,0)      {D1};
  \node[klein] (D2) at (0.6,-1.425) {D2};
  \node[klein] (D3) at (0.6,-3.325) {D3};
  \node[klein] (D4) at (0.6,-5.225) {D4};
  % ---------- Versuche und Abschnitte (nach Abschnitten gruppiert) ----------
  \node[vers] (VA) at (8.2,0)      {V1 bis V3};
  \node[vers, minimum height=16.5mm] (VB) at (8.2,-1.425) {V4 bis V7};
  \node[vers, minimum height=16.5mm] (VC) at (8.2,-3.325) {V8, V9};
  \node[vers, minimum height=16.5mm] (VD) at (8.2,-5.225) {V10, V11};
  \node[absch] (SA) at (11.3,0)      {Abschnitt~\ref{sec:eval_f1}};
  \node[absch, minimum height=16.5mm] (SB) at (11.3,-1.425) {Abschnitt~\ref{sec:eval_f2}};
  \node[absch, minimum height=16.5mm] (SC) at (11.3,-3.325) {Abschnitt~\ref{sec:eval_f3}};
  \node[absch, minimum height=16.5mm] (SD) at (11.3,-5.225) {Abschnitt~\ref{sec:eval_f4}};
  \node[ba block, minimum width=22mm, minimum height=64mm, text width=20mm] (AB) at (14.4,-2.85)
    {Abgleich\\Abschnitt\\\ref{sec:abgleich}};
  % ---------- Verbindungen ----------
  \foreach \a/\b in {D1/F1, D2/F2, D3/F3, D4/F4} \draw[ba pfeil] (\a) -- (\b);
  \draw[ba pfeil] (F1.east) -- (A2.west);
  \draw[ba pfeil] (F2.east) -- (A3.west);
  \draw[ba pfeil] (F2.east) -- (A4.west);
  \draw[ba pfeil] (F3.east) -- (A5.west);
  \draw[ba pfeil] (F3.east) -- (A6.west);
  \draw[ba pfeil] (F4.east) -- (A1.west);
  \draw[ba pfeil] (F4.east) -- (A7.west);
  \draw[ba pfeil] (A2.east) -- (A2 -| VA.west);
  \draw[ba pfeil] (A3.east) -- (A3 -| VB.west);
  \draw[ba pfeil] (A4.east) -- (A4 -| VB.west);
  \draw[ba pfeil] (A5.east) -- (A5 -| VC.west);
  \draw[ba pfeil] (A6.east) -- (A6 -| VC.west);
  \draw[ba pfeil] (A1.east) -- (A1 -| VD.west);
  \draw[ba pfeil] (A7.east) -- (A7 -| VD.west);
  \foreach \v/\s in {VA/SA, VB/SB, VC/SC, VD/SD} \draw[ba pfeil] (\v) -- (\s);
  \draw[ba pfeil] (SA.east) -- (SA -| AB.west);
  \draw[ba pfeil] (SB.east) -- (SB -| AB.west);
  \draw[ba pfeil] (SC.east) -- (SC -| AB.west);
  \draw[ba pfeil] (SD.east) -- (SD -| AB.west);
\end{tikzpicture}
 in einer figure-Umgebung. Nicht skalieren.
\begin{tikzpicture}[ba]
  \tikzset{
    klein/.style={ba block, minimum width=12mm, minimum height=7mm, inner sep=2pt},
    frage/.style={ba block hervor, minimum width=12mm, minimum height=7mm, inner sep=2pt},
    vers/.style={ba block, minimum width=21mm, inner sep=2pt},
    absch/.style={ba block, minimum width=21mm, inner xsep=5pt, inner ysep=2pt},
    kopf/.style={font=\footnotesize, text=baGrau, anchor=base},
  }
  % ---------- Spaltenköpfe ----------
  \node[kopf] at (0.6,0.65)  {Defizite};
  \node[kopf] at (2.9,0.65)  {Fragen};
  \node[kopf] at (5.2,0.65)  {Anforderungen};
  \node[kopf] at (8.2,0.65)  {Versuche};
  \node[kopf] at (11.3,0.65) {Abschnitte};
  \node[kopf] at (14.4,0.65) {Abgleich};
  % ---------- Anforderungen (Reihenfolge nach Fragen) ----------
  \node[klein] (A2) at (5.2,0)     {A2};
  \node[klein] (A3) at (5.2,-0.95) {A3};
  \node[klein] (A4) at (5.2,-1.9)  {A4};
  \node[klein] (A5) at (5.2,-2.85) {A5};
  \node[klein] (A6) at (5.2,-3.8)  {A6};
  \node[klein] (A1) at (5.2,-4.75) {A1};
  \node[klein] (A7) at (5.2,-5.7)  {A7};
  % ---------- Fragen und Defizite ----------
  \node[frage] (F1) at (2.9,0)      {F1};
  \node[frage] (F2) at (2.9,-1.425) {F2};
  \node[frage] (F3) at (2.9,-3.325) {F3};
  \node[frage] (F4) at (2.9,-5.225) {F4};
  \node[klein] (D1) at (0.6,0)      {D1};
  \node[klein] (D2) at (0.6,-1.425) {D2};
  \node[klein] (D3) at (0.6,-3.325) {D3};
  \node[klein] (D4) at (0.6,-5.225) {D4};
  % ---------- Versuche und Abschnitte (nach Abschnitten gruppiert) ----------
  \node[vers] (VA) at (8.2,0)      {V1 bis V3};
  \node[vers, minimum height=16.5mm] (VB) at (8.2,-1.425) {V4 bis V7};
  \node[vers, minimum height=16.5mm] (VC) at (8.2,-3.325) {V8, V9};
  \node[vers, minimum height=16.5mm] (VD) at (8.2,-5.225) {V10, V11};
  \node[absch] (SA) at (11.3,0)      {Abschnitt~\ref{sec:eval_f1}};
  \node[absch, minimum height=16.5mm] (SB) at (11.3,-1.425) {Abschnitt~\ref{sec:eval_f2}};
  \node[absch, minimum height=16.5mm] (SC) at (11.3,-3.325) {Abschnitt~\ref{sec:eval_f3}};
  \node[absch, minimum height=16.5mm] (SD) at (11.3,-5.225) {Abschnitt~\ref{sec:eval_f4}};
  \node[ba block, minimum width=22mm, minimum height=64mm, text width=20mm] (AB) at (14.4,-2.85)
    {Abgleich\\Abschnitt\\\ref{sec:abgleich}};
  % ---------- Verbindungen ----------
  \foreach \a/\b in {D1/F1, D2/F2, D3/F3, D4/F4} \draw[ba pfeil] (\a) -- (\b);
  \draw[ba pfeil] (F1.east) -- (A2.west);
  \draw[ba pfeil] (F2.east) -- (A3.west);
  \draw[ba pfeil] (F2.east) -- (A4.west);
  \draw[ba pfeil] (F3.east) -- (A5.west);
  \draw[ba pfeil] (F3.east) -- (A6.west);
  \draw[ba pfeil] (F4.east) -- (A1.west);
  \draw[ba pfeil] (F4.east) -- (A7.west);
  \draw[ba pfeil] (A2.east) -- (A2 -| VA.west);
  \draw[ba pfeil] (A3.east) -- (A3 -| VB.west);
  \draw[ba pfeil] (A4.east) -- (A4 -| VB.west);
  \draw[ba pfeil] (A5.east) -- (A5 -| VC.west);
  \draw[ba pfeil] (A6.east) -- (A6 -| VC.west);
  \draw[ba pfeil] (A1.east) -- (A1 -| VD.west);
  \draw[ba pfeil] (A7.east) -- (A7 -| VD.west);
  \foreach \v/\s in {VA/SA, VB/SB, VC/SC, VD/SD} \draw[ba pfeil] (\v) -- (\s);
  \draw[ba pfeil] (SA.east) -- (SA -| AB.west);
  \draw[ba pfeil] (SB.east) -- (SB -| AB.west);
  \draw[ba pfeil] (SC.east) -- (SC -| AB.west);
  \draw[ba pfeil] (SD.east) -- (SD -| AB.west);
\end{tikzpicture}
